\medskip

\textbf{Partie A.}

\medskip

\begin{enumerate}
\item On a $u_1 - u_0 = u_2 - u_1 = u_3 - u_2 = 25$ : ce nombre 25 est la raison de la suite arithmétique $(u_n)$.
\item Dans 2 heures il y a $12 \times 10$ minutes : $u_0$ est donc augmenté de $12 \times 25 = 300$ champignons. Il y aura donc au bout de 2 h : $100 + 300 = 400 = 4\times 100 = 4 \times u_0$ : la population aura donc quadruplé en 2~heures.
\end{enumerate}

\medskip

\textbf{Partie B.}

\medskip

\begin{enumerate}
\item On a $v_0 = 100$, $v_1 = 2 \times u_0, \, v_2 = 2v_1, \, v_3 = 2v_2$ : ces égalités montrent que les termes $v_0,\: v_1,\: v_2,\: v_3$ sont en progression géométrique de raison 2.
\item Seul le graphique 1 est susceptible de représenter la suite $(v_n)$.
\item 4~h représentent 240~min soit $6 \times 40$~min ; il y aura donc au bout de 4~h,

$u_6 = 100 \times 2^6 = 100 \times 64 = \np{6400}$~champignons.
\item Après 4~h 40~min la population sera le double de celle de 4~h, soit : $2 \times  \np{6400} = \np{12800}$.

40 min plus tard, soit après 5~h 20~min la population aura encore doublé et sera égale à $2 \times \np{12800} = \np{25600}$.

Comme $\np{12800} < \np{18000} < \np{25600}$, cette population de \np{18000} champignons est cohérente avec l'hypothèse du doublement de la population toutes les quarante minutes.
\end{enumerate}

\bigskip
